% Image credits for the photographs and AI illustrations of Book 4
% (University Biology, Year 2). Machine-readable license records live in
% images/book4/CREDITS.md; keep the two files in sync.
\chapter*{Créditos das imagens}

% \sloppy must be in force when the paragraph is BROKEN, hence the \par before
% \endgroup -- without it the group closes first and the tolerance is restored.
\begingroup\sloppy

Os desenhos deste livro são originais. As fotografias são usadas sob
as respectivas licenças livres, com agradecimento aos seus autores:

\begin{itemize}
  \item Capítulo~\ref{ch:b2:microbial-metabolism}: fotografia de Sergei
        Winogradsky, década de 1890 --- domínio público (Wikimedia Commons).
  \item Capítulo~\ref{ch:b2:genome-diversification}: fotografia de
        Barbara McClintock em seu laboratório, Smithsonian Institution /
        Science Service, 1947 --- domínio público (Wikimedia Commons).
  \item Capítulo~\ref{ch:b2:blood-circulation}: retrato de William
        Harvey e a prancha das veias do antebraço de sua
        \emph{Exercitatio anatomica de motu cordis} (1628), Wellcome
        Collection --- CC~BY~4.0 (Wikimedia Commons); micrografia
        eletrônica de varredura de uma hemácia, uma plaqueta e um
        linfócito, Electron Microscopy Facility, National Cancer
        Institute --- domínio público (Wikimedia Commons).
  \item Capítulo~\ref{ch:b2:speciation}: prancha dos tentilhões de Darwin
        de John Gould, do \emph{Journal of Researches} (1845) --- domínio
        público (Wikimedia Commons).
  \item Capítulo~\ref{ch:b2:biogeochemical-cycles}: o Observatório de
        Mauna Loa (Forrest M. Mims III, 2006) e uma amostra em frasco
        sendo coletada ali (comandante John Bortniak, 1982), NOAA Photo
        Library --- domínio público (Wikimedia Commons).
\end{itemize}

Os links completos das fontes e os endereços das licenças de cada
fotografia estão em \texttt{images/\allowbreak book4/\allowbreak CREDITS.md},
no repositório do livro.

\bigskip
Além das fotografias e dos desenhos originais, cerca de setenta
figuras espalhadas pelos capítulos do volume são ilustrações geradas
por inteligência artificial, produzidas com a geração de imagens da
OpenAI a partir de instruções escritas pelos editores do livro, e
revisadas quanto à exatidão biológica antes da inclusão. A instrução
completa de cada imagem gerada está registrada em
\texttt{images/\allowbreak book4/\allowbreak ai/\allowbreak PROMPTS.md}, no repositório.
\par\endgroup
\clearpage
